\subsection{局部凸空间中凸集的分离}

命题 4. --- 设 A 是局部凸空间 E 中的一个非空闭凸集，K 是 E 中一个不与 A 相交的非空紧凸集。则存在一个把 A 与 K 严格分离的闭超平面 H。

因为存在 E 中 0 的一个开凸邻域 V，使 A + V 与 K + V 不相交 (GT, II, § 4.3, prop. 4)。由于 A + V 与 K + V 在 E 中是凸的且开的，II, 第 37 页的 prop. 1 表明存在一个把 A + V 与 K + V 严格分离的闭超平面 H，从而更是把 A 与 K 严格分离。

注记. --- 在 Hausdorff 局部凸空间 E 中，设 A 和 B 是两个不相交的非空闭凸集，若 E 是有限维的，则存在一个把 A 与 B 分离的闭超平面 (II, 第 78 页, exerc. 13)；但当 E 是无限维时，这个结论未必成立 (II, 第 78 页, exerc. 10 及 11)。

推论 1. --- 在局部凸空间中，每个闭凸集 A 是包含它的诸闭半空间的交。

事实上，对每一点 \(x \notin \mathbf{A}\) ，存在一个把 \(x\) 与 \(\mathbf{A}\) 严格分离的闭超平面（用 prop. 4）。

推论 2. --- 在 Hausdorff 局部凸空间中，每个紧凸集 A 是包含它且由 A 的支撑超平面确定的诸闭半空间的交。

因为，设 \(x_{0} \notin A\) ；\(\{x_{0}\}\) 是闭的，因此存在一个把 \(x_{0}\) 与 A 严格分离的闭超平面 H (prop. 4)；设 \(f(x) = \alpha\) 是 H 的一个方程（f 为连续线性型），并设 \(f(x) > \alpha\) （对所有 \(x \in A\) ）。若令 \(\gamma = \inf_{x \in A} f(x)\) ，则由 \(f(x) \geqslant \gamma\) 定义的半空间包含 A，由方程为 \(f(x) = \gamma\) 的支撑超平面确定，且不含 \(x_{0}\) ；推论得证。

在局部凸空间中，一个既非紧又无内点的闭凸集可能没有任何闭的支撑超平面 (II, 第 86 页, exerc. 18：另参见 V, 第 71 页, exerc. 11)。

推论 3. --- 在局部凸空间中，每个线性流形 M 的闭包是包含 M 的诸闭超平面的交。

对所有 \(x \notin \overline{\mathbf{M}}\) ，设 \(\mathbf{H}\) 是一个把 \(x\) 与 \(\overline{\mathbf{M}}\) 严格分离的闭超平面；

于是 \(\overline{M}\) 平行于 H；闭超平面 \(H_{1}\) （包含 \(\overline{M}\) 且平行于 H 的）不含 x。推论得证。

推论 4. --- 设 C 是局部凸空间 E 中的一个闭凸集。E 的子集 A 含于 C，必须且只需对 E 上每个使得对所有 C 中的 x 有 \(u(x) \geqslant 0\) 的实值连续仿射函数 u，都有对所有 A 中的 y，\(u(y) \geqslant 0\) 。

条件显然是必要的。反之我们证明它是充分的；若一点 \(x \in A\) 不含于 C，则存在一个方程为 \(f(z) = \alpha\) 的闭超平面把 x 与 C 严格分离；若例如设 \(f(x) < \alpha\) ，则连续仿射函数 \(u = f - \alpha\) 与假设矛盾。

推论 5. --- 在局部凸空间 E 中，每个以 0 为顶点的凸锥 C 的闭包是包含 C 且由过 0 的闭超平面确定的诸闭半空间的交。

因为 \(\overline{\mathbf{C}}\) 是一个以 0 为顶点的凸锥 (II, 第 13 页, prop. 14)。对 \(x \notin \overline{\mathbf{C}}\) ，存在一个把 \(x\) 与 \(\overline{\mathbf{C}}\) 严格分离的闭超平面 H (prop. 4)。这时只须应用下面的引理：

引理 1. --- 若一个以 0 为顶点的锥 A 含于由超平面 H 确定的一个开半空间中，则它含于由一个平行于 H 且过 0 的超平面 \(H_{0}\) 确定的闭半空间中。

设 \(f(z) = \alpha\) （\(\alpha < 0\) ）是 H 的一个方程，于是 \(f(z) = 0\) 是 \(H_0\) 的方程。若存在 \(z \in A\) 使 \(f(z) < 0\) ，则将存在 \(\lambda > 0\) 使 \(f(\lambda z) = \alpha\) ，而由于 \(\lambda z \in A\) ，这将与假设矛盾。

